
\begin{abstract}
Cloud-hosted ML services enable access to powerful models through remote APIs, but leave clients unable to directly verify the computations that produce their predictions. Zero-knowledge machine learning (zkML) addresses this problem by allowing the provider to return each prediction with a cryptographic proof of correct execution. Proof generation, however, remains computationally expensive and memory intensive. Many zkML pipelines prove an entire inference in a single end-to-end invocation, constraining each request to the runtime and memory capacity of one machine.

We present \TITLE{}, a distributed proof-serving system that replaces a full-model proof with smaller proving jobs, each covering part of the model computation. \TITLE{} materializes the intermediate tensors needed at partition boundaries before proving begins, allowing jobs to execute independently across prover machines. Cryptographic commitments bind adjacent partitions so that the resulting proofs establish one consistent end-to-end inference. \TITLE{} also reuses compiled circuits and proving keys across requests and schedules jobs using runtime and memory estimates.

Our characterization shows that partitioning reduces aggregate prover time and peak memory even before parallel execution. Across convolutional, generative, and transformer workloads, \TITLE{} reduces end-to-end latency by up to \(5.6\times\) on a single worker. Distributing jobs across multiple workers yields up to \(69\times\) lower latency than monolithic proving and reduces peak per-machine memory by up to \(48\times\). These results show that scalable zkML proof generation depends not only on the proving backend, but also on how computation is partitioned and executed across machines.
\end{abstract}
